% =====================================================================
\section{Numerical Experiments}
\label{sec:exp}

\subsection{Scope of the Results}
\label{sec:exp-scope}
The results come from a 30-day continuous diagnostic run with random seed
9{,}900{,}950, which was not used in training. We compare the middle 28
days and use the first and last days only to connect the continuous
boundaries. During evaluation, weights are frozen and exploration is set
to zero. All policies share the same truck, vessel, and service random
streams on a given day. We record the four cost terms, demand level,
transaction types, and unfinished jobs for each day. The source code,
raw policy results, and figure scripts are released with the paper.

\subsection{Confirmation of Architectural Operation}
\label{sec:exp-confirm}
Implementation checks, commitment constraints, and automated tests
confirm four structural properties. Proposal and acceptance use separate
models. One-sided consent never produces a commitment. Realised future
arrival and completion times are excluded from policy inputs. State also
continues across day boundaries. The run logs confirm two additional
properties. Frozen-policy evaluation made \textbf{zero} counterfactual
simulation calls. The identity check also passed at all four sampled
points with counterfactual labels: a branch that changed no action
reproduced the action and cost of the main run.

\subsection{Cost Reduction and Paired Daily Comparison}
\label{sec:exp-paired}
Throughout this section, \emph{learned policy} denotes the full
architecture of Sect.~\ref{sec:arch}---proposal, acceptance, and central
commitment---and not the proposal policy, which is one of its two
networks. Table~\ref{tab:policies} compares each policy's 28-day total
cost with no reallocation. A positive value means a cost reduction, whereas a
negative value means a cost increase. The rows are grouped by research
question rather than ranked in one list.

\begin{center}
\begin{minipage}{\textwidth}
\centering
\captionof{table}{Total cost reduction by policy. Each reduction is measured
against the 28-day total cost without reallocation. Positive values mean
lower cost, whereas negative values mean higher cost.}
\label{tab:policies}
\small
\begin{tabular}{@{}llr@{}}
\toprule
Policy & Action range & Reduction vs. no reallocation \\
\midrule
\textbf{Learned policy} & block\,+\,time & \textbf{+14.72\%} \\
Learned, time only & time & +13.44\% \\
Learned, block only & block & $-$1.31\% \\
\midrule
Untrained model & block\,+\,time & +7.97\% \\
\midrule
Rule: least-loaded slot & time & +3.41\% \\
Rule: least-loaded block\,+\,slot & block\,+\,time & +1.51\% \\
Rule: least slack & block & +0.53\% \\
Rule: first-come-first-served & block & +0.17\% \\
Rule: shortest processing time & block & +0.05\% \\
Rule: net gain & block & $-$0.48\% \\
\midrule
No reallocation & --- & 0.00\% \\
\bottomrule
\end{tabular}
\end{minipage}
\end{center}
Over 28 days, total cost was KRW 10.00\,bn without reallocation and KRW
8.52\,bn with the learned policy. The reduction was KRW 1.471\,bn, or
14.7\%. Time adjustment alone reduces cost by 13.44\%, whereas block
reallocation alone raises cost by 1.31\%. The total reduction therefore
comes mainly from arrival-time adjustment.
We run every policy with the same random streams to reduce simulation
noise, following Kleijnen's common-random-number method~\cite{ref22}.
Days are treated as independent comparison units. We report two paired
tests because they answer different questions. A two-sided sign test
counts only the number of days on which one policy is cheaper, so it is
blind to how large each day's difference is. A Wilcoxon signed-rank test
ranks the differences by magnitude before summing them, so a policy that
wins on fewer days but by a wider margin is scored differently by the
two. The distinction matters here because the learned policy acts
selectively: it is not designed to win on every day but to reduce cost
sharply on congested ones. With 28 days and no tied differences, both
tests are computed exactly. Table~\ref{tab:paired} reports both.

\begin{center}
\begin{minipage}{\textwidth}
\centering
\captionof{table}{Paired daily comparison between the learned policy and each
comparison policy. ``Learned lower'' counts the days on which the
learned policy costs less. Both $p$ values are two-sided and exact: a
sign test, which uses only the direction of each daily difference, and a
Wilcoxon signed-rank test, which also uses its magnitude. Values below
0.05 are in bold.}
\label{tab:paired}
\small
\setlength{\tabcolsep}{4pt}
\begin{tabular}{@{}>{\raggedright\arraybackslash}p{0.185\textwidth}>{\raggedright\arraybackslash}p{0.30\textwidth}rrr@{}}
\toprule
Purpose & Comparison policy & \shortstack{Learned lower\\days} &
\shortstack{Sign test\\$p$} & \shortstack{Signed-rank\\$p$} \\
\midrule
Learning effect & Untrained model & 18/28 & 0.185 & 0.164 \\
Same action range & Least-loaded block\,+\,slot rule & 20/28 & \textbf{0.036} & \textbf{0.015} \\
Block rule & Least slack & 20/28 & \textbf{0.036} & \textbf{$<$0.001} \\
Block rule & First-come-first-served & 22/28 & \textbf{0.004} & \textbf{$<$0.001} \\
Block rule & Shortest processing time & 21/28 & \textbf{0.013} & \textbf{0.006} \\
Block rule & Net gain & 21/28 & \textbf{0.013} & \textbf{0.003} \\
Baseline & No reallocation & 26/28 & \textbf{$<$0.001} & \textbf{$<$0.001} \\
\bottomrule
\end{tabular}
\end{minipage}
\end{center}
The learned policy costs less than the rule with the same block-and-time
action range on 20 of 28 days (sign test $p=0.036$, signed-rank
$p=0.015$). It costs less than the untrained model on 18 days, but this
difference is not statistically clear under either test ($p=0.185$ and
$p=0.164$). At the 5\% level the two tests agree on every row of
Table~\ref{tab:paired}. They are not equally sensitive, however: the
signed-rank $p$ is the smaller of the two in six of the seven rows, which
is what a concentrated effect produces. The architecture's
action range and the additional effect of learning must therefore be
interpreted separately. The learned policy is more expensive on only
two days, and both differences are small, whereas the largest reductions
all occur on congested days (Fig.~\ref{fig:paired}).

\begin{center}
\begin{minipage}{\textwidth}
\centering
\includegraphics[width=\textwidth]{fig-paired-daily.pdf}
\captionof{figure}{Paired daily cost difference against no reallocation over
the same 28 days with the same random streams. Days are sorted by the
difference, and a positive value means the learned policy costs less.
(a) All 28 days; four congested days dominate the axis, so (b) enlarges
the shaded ranks 1--24. The learned policy costs less on 26 of the 28
days (two-sided sign test, $p<0.001$). The two exceptions are small,
KRW $-17.1$\,m and KRW $-0.5$\,m, whereas each of the four congested days
exceeds KRW 290\,m.}
\label{fig:paired}
\end{minipage}
\end{center}

\subsection{Demand Level and the Decomposition of Actions}
\label{sec:exp-decomp}
Table~\ref{tab:actions} decomposes the reduction by demand level. The
four days with 12{,}500 or 15{,}000 trucks account for 90.5\% of the
total. This concentration supports developing the architecture toward a
policy that \emph{intervenes selectively when congestion is observed}
rather than continuously. Restricting the policy to one action type at a
time separates the spatial and temporal contributions. Arrival-time
adjustment accounts for 91.3\% of the total reduction, mostly on the four
congested days (Fig.~\ref{fig:decomp}). This share should be read
together with the cost structure of Sect.~\ref{sec:env-cost}, where truck
waiting accounts for 96.4\% of total cost. An action that moves waiting
time therefore acts on the dominant term, and the relative standing of
spatial and temporal actions may differ under a different cost mix. Block-only reallocation saves KRW 0.124\,bn under smooth
and normal demand but raises cost at the highest demand, producing a net
increase of KRW 0.131\,bn. The policy should therefore adjust its use of
spatial and temporal actions by demand level. Time-slot quotas were open
in this diagnostic run. The arrival-time results thus describe the
\emph{potential range} of time shifting, not an executable schedule.
Sensitivity analysis with finite quotas can define the practical range.

\begin{table}[tb]
\centering
\caption{Cost reduction by demand level (KRW bn, relative to no
reallocation). Columns three to five separate the allowed actions. The
last column compares the trained and untrained models. Arrival-time
adjustment and learning contribute mainly on the four congested days,
whereas block reallocation raises cost at the highest demand.}
\label{tab:actions}
\small
\begin{tabular}{@{}rrrrrr@{}}
\toprule
Daily & & Both & Time & Block & Due to \\
demand & Days & actions & only & only & training \\
\midrule
3,500 & 12 & 0.047 & $-$0.013 & 0.057 & 0.056 \\
5,000 & 8 & 0.060 & 0.029 & 0.067 & 0.029 \\
7,500 & 4 & 0.032 & 0.008 & 0.000 & $-$0.049 \\
12,500 & 2 & 0.611 & 0.573 & $-$0.014 & 0.144 \\
15,000 & 2 & 0.722 & 0.747 & $-$0.241 & 0.495 \\
\midrule
\textbf{Total} & \textbf{28} & \textbf{1.471} & \textbf{1.344} & \textbf{$-$0.131} & \textbf{0.675} \\
\multicolumn{2}{@{}l}{\textit{share of total}} & 100\% & \textbf{91.3\%} & --- & \textbf{45.9\%} \\
\bottomrule
\end{tabular}
\end{table}

\begin{center}
\begin{minipage}{\textwidth}
\centering
\includegraphics[width=\textwidth]{fig-decomposition.pdf}
\captionof{figure}{Cost reduction by demand level, decomposed by action
range. The values are those of Table~\ref{tab:actions}, in million KRW
(1\,bn $=$ 1{,}000\,m). All five demand levels share one axis, and the
shaded band marks the two congested levels. Because those two are more
than ten times larger than the rest, the value is printed next to each
short bar. The number of days observed appears as $n$ under each tick.
Bars measure a
restricted policy against no reallocation. The diamond is \emph{not} a
fourth action: it is the trained model minus the untrained one---the
training contribution, that is, the last column of
Table~\ref{tab:actions}---and carries a different mark for that reason.
Moving from the non-congested to the congested regime, arrival-time
adjustment and the training contribution both grow sharply, whereas
block-only reallocation turns negative.}
\label{fig:decomp}
\end{minipage}
\end{center}

\subsection{Rule-Based Policies of the Same Action Range}
\label{sec:exp-rules}
To determine whether the difference in Sect.~\ref{sec:exp-decomp} comes
from the action itself, the rule-based policies must receive the same
action range. We therefore transfer the \emph{least-loaded} criterion
from the spatial axis to the temporal axis and add two rules. This change
introduces no new parameter and uses the same congestion trigger as the
other rules. Three findings follow. First, arrival-time adjustment under
a rule reduces cost by 3.41\% relative to no reallocation
($p=0.0125$). The temporal action therefore has value by itself. Second,
with the full spatio-temporal action range matched, the learned policy
costs less than the rule on 20 of 28 days ($p=0.036$). Third, for
arrival-time adjustment alone, the rule costs less on more days; the
learned policy is lower on only 8 of 28 days. Yet the learned policy's
total reduction is about four times larger. The rule is slightly better
under smooth demand, whereas the learned policy reduces cost 3.8--5.6
times more on congested days. Thus, the learned policy tends to produce
large savings on expensive days rather than frequent small wins. Both
the number of wins and the total reduction are needed to show this
pattern. The rule using both actions also weakens under congestion. This
suggests that the congestion sensitivity of block reallocation may be a
property of the action, not only of one policy.

\subsection{Learning Effect and Dispatching Sensitivity}
\label{sec:exp-learning}
We first examine model fit. Across 30 training iterations, the proposal
network's Huber loss decreases from 0.079 to 0.029 on training decisions
and from 0.268 to 0.119 on validation decisions. The acceptance
network's loss decreases from 0.171 to 0.027 and from 0.315 to 0.136,
respectively. These values are means over the first and last five
iterations, with targets in units of KRW 100{,}000. The decline in
validation loss indicates that the fit extends to unseen decisions.
However, final validation loss remains about four times training loss.
With about 56 labelled decisions per iteration, this gap supports the
choice of a small model. As exploration decreases, the share of
decisions with exactly zero counterfactual difference falls from 51.4\%
to 37.6\%. Training runs 4{,}671 counterfactual worlds in total.
Fig.~\ref{fig:learning} shows both loss curves together with the
exploration schedule that accompanies them.

\begin{center}
\begin{minipage}{\textwidth}
\centering
\includegraphics[width=\textwidth]{fig-learning-curve.pdf}
\captionof{figure}{Training and validation loss over 30 iterations, with the
per-iteration value drawn faintly beneath a five-iteration mean. (a) The
proposal network and (b) the acceptance network are drawn separately
because they are fitted on different decisions, and the vertical axis is
logarithmic because the loss spans three orders of magnitude. (c) The
exploration probability $\varepsilon$ decays linearly from 0.50 to 0.05.
Shading marks the iterations whose demand level was drawn at 12{,}500
trucks per day or above. Because the demand level is redrawn every
iteration, much of the visible fluctuation, including the rise around
iterations 17--24, follows that draw rather than the fit.}
\label{fig:learning}
\end{minipage}
\end{center}
Cost is KRW 10.00\,bn without reallocation, KRW 9.20\,bn with the
untrained model, and KRW 8.52\,bn with the trained model. Training
therefore adds KRW 0.675\,bn, or 45.9\% of the total reduction. However,
the trained model costs less on only 18 of 28 days (sign test
$p=0.185$, signed-rank $p=0.164$). The
mean daily difference is KRW $-24.1$\,m, about 18 times the median of
$-1.3$\,m. This gap shows that the learning effect is concentrated rather
than evenly distributed. The four congested days account for 94.6\% of
the learning effect, and the trained model costs less on all four. The
corresponding counts at lower demand levels are 9/12, 5/8, and 0/4. For
the remaining 24 days, the mean difference is $-1.5$\,m and the
bootstrap interval includes zero.
A stratified bootstrap, which resamples the daily differences within
demand strata with 20{,}000 resamples, gives a 95\% percentile interval
of $[-32,-17]$\,m per day for the mean difference. Stratification
preserves the demand mix fixed by the design. Without it, the interval
widens to $[-52,-3]$\,m, and 99.3\% of resamples remain negative. Both intervals exclude zero, but the stratified interval is
partly narrow by construction because each congested stratum contains
only two days. The interpretation is therefore limited. Learning
produces a large total reduction and is lower on all four congested
days, but the day-level sign test is not significant. Establishing the
learning effect at a conventional significance level requires more
congested days, not simply more days overall.
The ranking of crane-dispatching rules also changes with demand. The
shortest processing time rule is 22--25\% more expensive than the other
rules under smooth and normal demand, but it is the least expensive from
7{,}500 trucks upward. It also has the lowest overall cost under the
designed demand mix. The effect of reallocation should therefore be
tested under other dispatching rules.
